
\uline{entier} @@M0@@ \uline{tel que pour tout} @@M1@@, \uline{il existe une section} @@M2@@ \uline{qui soit une base de} @@M3@@ \uline{pour tout} @@M4@@.

On peut trouver un entier @@M5@@, et une section @@M6@@ telle que pour tout @@M7@@, @@M8@@, et que @@M9@@ soit un ouvert affine (EGA II 4.5.4). Sur @@M10@@, on peut trouver une section @@M11@@ satisfaisant les conditions du lemme 6.7; de plus, il existe un entier @@M12@@ tel que la section @@M13@@ se prolonge à @@M14@@ tout entier. Posons @@M15@@; alors @@M16@@; en faisant le produit tensoriel de la section de @@M17@@ obtenue plus haut et de @@M18@@, on obtient donc une section de @@M19@@ pour tout @@M20@@. Comme pour tout @@M21@@, @@M22@@, et que la section @@M23@@ est une base de @@M24@@ pour tout @@M25@@, les sections de @@M26@@ ainsi construites pour tout @@M27@@ forment une base de @@M28@@ pour tout @@M29@@.

\textbf{Théorème 6.9.} \uline{Soit} @@M30@@ \uline{un schéma noethérien de dimension} @@M31@@, \uline{ayant un faisceau ample} @@M32@@. \uline{Alors} @@M33@@, \uline{et à fortiori} @@M34@@.

(Cette dernière relation est due à J.P. Serre, dans le cas général envisagé ici).

Nous allons d'abord montrer que si @@M35@@ sont deux faisceaux inversibles sur @@M36@@, de classes @@M37@@, alors @@M38@@.

Soient @@M39@@ un fermé de @@M40@@, @@M41@@ ses composantes irréductibles, et pour tout @@M42@@, soit @@M43@@ un point de @@M44@@. D'après le lemme 6.8, il existe des entiers @@M45@@, tels que, quels que soient les entiers @@M46@@, on puisse trouver des homomorphismes de @@M47@@ dans @@M48@@, @@M49@@ qui soient des isomorphismes aux points @@M50@@; on obtient donc des morphismes:
@@M51@@
